\documentclass[letterpaper]{article}

\usepackage{natbib,alifeconf}
\usepackage{graphicx,amsmath,amssymb,booktabs,xcolor,multirow}
\usepackage{url,hyperref}

% Tighten float spacing
\setlength{\textfloatsep}{6pt plus 2pt minus 2pt}
\setlength{\floatsep}{6pt plus 2pt minus 2pt}
\setlength{\intextsep}{6pt plus 2pt minus 2pt}

% Conditional figure include (shared with main)
\newcommand{\figifexists}[4]{%
  \IfFileExists{#1}{%
    \includegraphics[width=#2,#3]{#1}%
  }{%
    \begin{minipage}{#2}\vspace{2em}\centering\footnotesize[Fig.\ pending: #4]\vspace{2em}\end{minipage}%
  }%
}

\title{Supplementary Material: Semi-Life: A Capability Ladder\\
for the Virus-to-Life Transition}

\author{Anonymous}

\begin{document}
\InputIfFileExists{auto_macros}{}{}
\maketitle

\noindent
This supplement accompanies the main paper and contains supporting details
referenced therein.
Table~\ref{tab:lookup} maps main-paper claims to supplement sections.

\begin{table}[h]
\centering
\footnotesize
\caption{Lookup table: main paper reference $\to$ supplement section.}
\label{tab:lookup}
\begin{tabular}{ll}
\toprule
Main paper reference & Supplement section \\
\midrule
Archetype parameters & S1 \\
Extended hypothesis rationales \& results & S2 \\
Shock resilience analysis & S3 \\
Effect size robustness check & S4 \\
Pre-registered sham controls & S5 \\
V4 policy evolution & S6 \\
Inter-entity competition & S7 \\
V1 protection regime sweep & S8 \\
\bottomrule
\end{tabular}
\end{table}

% ============================================================================
\section*{S1. Archetype Parameter Table}
\label{sec:s1}

\begin{table}[h]
\centering
\footnotesize
\caption{Archetype parameter summary. Rows labelled ``shared'' apply to all archetypes.}
\label{tab:archetypes-s}
\resizebox{\columnwidth}{!}{%
\begin{tabular}{lccc}
\toprule
Parameter & Viroid & Virus & ProtoOrganelle \\
\midrule
Baseline capabilities & V0 & V0+V1 & V1+V2+V3 \\
\texttt{replication\_threshold} & 0.60 & 0.60 & 0.80 \\
\texttt{replication\_cost} & 0.27 & 0.27 & 0.30 \\
\texttt{boundary\_decay\_rate} & 0.002 & 0.001 & 0.002 \\
\texttt{boundary\_repair\_rate} & 0.010 & 0.010 & 0.010 \\
\midrule
\multicolumn{4}{l}{\emph{V1 protective parameters (shared)}} \\
\texttt{energy\_leakage\_rate} & \multicolumn{3}{c}{0.005} \\
\texttt{env\_damage\_probability} & \multicolumn{3}{c}{0.05} \\
\texttt{env\_damage\_amount} & \multicolumn{3}{c}{0.05} \\
\texttt{boundary\_damage\_absorption} & \multicolumn{3}{c}{0.8} \\
\texttt{boundary\_damage\_integrity\_cost} & \multicolumn{3}{c}{0.02} \\
\midrule
\multicolumn{4}{l}{\emph{V2 overconsumption parameters (shared)}} \\
\texttt{overconsumption\_waste\_fraction} & \multicolumn{3}{c}{0.3} \\
\texttt{optimal\_uptake\_rate} & \multicolumn{3}{c}{0.015} \\
\midrule
\multicolumn{4}{l}{\emph{V4 chemotaxis parameters (shared)}} \\
\texttt{v4\_move\_cost} & \multicolumn{3}{c}{0.01} \\
\texttt{v4\_max\_speed} & \multicolumn{3}{c}{1.0} \\
\texttt{v4\_mutation\_sigma} & \multicolumn{3}{c}{0.05} \\
\midrule
\multicolumn{4}{l}{\emph{V5 lifecycle parameters (shared)}} \\
\texttt{v5\_activation\_threshold} & \multicolumn{3}{c}{0.6} \\
\texttt{v5\_dispersal\_age} & \multicolumn{3}{c}{100} \\
\texttt{v5\_dispersal\_duration} & \multicolumn{3}{c}{20} \\
\texttt{v5\_dormant\_decay\_mult} & \multicolumn{3}{c}{0.3} \\
\texttt{v5\_dispersal\_decay\_mult} & \multicolumn{3}{c}{1.5} \\
\texttt{v5\_dispersal\_speed\_mult} & \multicolumn{3}{c}{2.0} \\
\bottomrule
\end{tabular}%
}
\end{table}

Key differentiating parameters across archetypes: Viroid and Virus share a
replication threshold of 0.60, while ProtoOrganelle uses 0.80 (reflecting
its more complex metabolic machinery).
Virus has a slower boundary decay rate (0.001 vs.\ 0.002) reflecting greater
capsid stability.
All shared parameters (V1 protective, V2 overconsumption, V4 chemotaxis,
V5 lifecycle) are identical across archetypes so that inter-archetype
comparisons isolate the effect of baseline capability composition.

% ============================================================================
\section*{S2. Extended Hypothesis Rationales and Results}
\label{sec:s2}

This section provides the full rationale paragraphs for each hypothesis
and per-harshness narrative detail that was condensed in the main paper.

\paragraph{H1 rationale (boundary cost--benefit tradeoff).}
V1 boundary repair costs per-step energy, but V1 also protects against
energy leakage and environmental damage.
In resource-rich environments where damage pressure is low relative to
repair costs, V1 is a net cost; in harsh environments, V1's protective
benefit dominates.

\paragraph{H1 extended results.}
In rich environments, V1 significantly reduces survival
($\delta=\statDeltaHoneRich{}$, $p_\text{corr}\statPHoneRich{}$):
mean alive drops from \aliveVzeroRich{} (V0) to \aliveVzeroVoneRich{}
(V0+V1), because boundary maintenance cost exceeds the leakage/damage
protection benefit when resources are abundant.
In medium, the effect is weaker but still significant
($\delta=\statDeltaHoneMedium{}$, $p_\text{corr}=\statPHoneMedium{}$):
\aliveVzeroMedium{}~$\to$~\aliveVzeroVoneMedium{}.
In sparse and scarce, V1 has \emph{no significant effect}
($\delta=\statDeltaHoneSparse{}$, $p=\statPHoneSparse{}$): boundary
protection exactly offsets its cost, yielding near-identical populations
($\sim$10.0 alive).

\paragraph{H2 rationale (metabolism buffering).}
V3 internal pool decouples replication threshold from instantaneous uptake,
buffering against resource fluctuations.

\paragraph{H2 extended results.}
V3 addition transforms survival: in rich environments, mean alive rises from
\aliveVzeroVoneVtwoRich{} (V0+V1+V2) to \aliveVzeroVoneVtwoVthreeRich{}
(V0+V1+V2+V3). In medium, the jump is
\aliveVzeroVoneVtwoMedium{}~$\to$~\aliveVzeroVoneVtwoVthreeMedium{}.

\paragraph{H3 rationale (replication liberation).}
Without V0, ProtoOrganelle cannot replicate regardless of metabolic state;
V0 addition activates replication capability already supported by the
existing V1--V3 infrastructure.

\paragraph{H3 extended results.}
Confirmed in rich ($\delta=1.00$, $p_\text{corr}<10^{-37}$) and medium
($\delta=1.00$, $p_\text{corr}<10^{-38}$), where liberated ProtoOrganelle
reaches \aliveProtoLibRich{} and \aliveProtoLibMedium{} mean alive respectively.
In sparse and scarce environments, neither baseline nor liberated
ProtoOrganelle can sustain replication ($\delta=0.00$, $p=1.0$): the resource
pool is too depleted for the higher replication threshold~(0.80).

\paragraph{H4 rationale (monotonic trend V0--V3).}
The Jonckheere-Terpstra test detects a dominant upward trend; pairwise
reversals at intermediate V-levels (e.g., V0$\to$V0+V1 dip) are present
and discussed in the main paper.

\paragraph{H5 rationale (chemotaxis benefit).}
V4 gradient-following enables directed movement toward resource-rich areas,
improving energy intake in heterogeneous environments.

\paragraph{H5 extended results.}
V4 chemotaxis improves survival where resources are heterogeneous:
mean alive rises from \aliveVzeroVoneVtwoVthreeMedium{} to
\aliveVzeroToVfourMedium{} in medium and from \aliveVzeroVoneVtwoVthreeSparse{}
to \aliveVzeroToVfourSparse{} in sparse.
However, in resource-rich environments where resources are uniformly abundant,
the per-step movement cost ($0.01 \times |\Delta\mathbf{v}|$) reduces survival
from \aliveVzeroVoneVtwoVthreeRich{} to \aliveVzeroToVfourRich{}.

\paragraph{H6 rationale (lifecycle staging).}
V5 dormancy conserves energy during scarcity, and dispersal enables
colonisation of resource patches, both improving population persistence.

\paragraph{H6 extended results.}
Confirmed only in rich ($\delta=1.00$, $p_\text{corr}<10^{-33}$), where V5
dormancy enables energy conservation during transient low-resource periods,
pushing mean alive from \aliveVzeroToVfourRich{} to \aliveVzeroToVfiveRich{}.
In medium ($\delta=\statDeltaHsixMedium{}$), sparse
($\delta=\statDeltaHsixSparse{}$), and scarce
($\delta=\statDeltaHsixScarce{}$), V5 \emph{reduces} survival: dormancy
phases suppress replication windows, and dispersal's elevated energy decay
($\times 1.5$) is counterproductive.

\paragraph{H7 rationale (full monotonic trend V0--V5).}
``Monotonic'' refers to the dominant aggregate trend detected by JT;
pairwise reversals at V4 and V5 are present.

\paragraph{H8 rationale (V2 overconsumption regulation).}
V2's regulator reduces overconsumption waste, providing a survival advantage
independent of V3 metabolism.

\paragraph{H8 extended results.}
Confirmed in rich ($\delta=\statDeltaHeightRich{}$,
$p_\text{corr}\statPHeightRich{}$): mean alive rises from
\aliveVzeroVoneRich{} (V0+V1) to \aliveVzeroVoneVtwoRich{} (V0+V1+V2).
In medium, sparse, and scarce, V2 has no significant effect
($\delta \approx 0$, $p=1.0$): overconsumption rarely occurs when
resources are scarce.

% ============================================================================
\section*{S3. Shock Resilience Analysis}
\label{sec:s3}

\begin{figure*}[t]
\centering
\figifexists{figures/fig_semi_life_recovery.pdf}{\textwidth}{}{shock recovery}
\caption{Recovery under periodic resource shocks (Viroid, sparse harshness,
  $n\!=\!100$).
  Left: shock period~$=50$ (rapid crash--recovery cycles).
  Centre: shock period~$=200$ (slow cycles).
  Right: no-shock baseline (same harshness level for comparison).
  Each line is one capability level; Okabe-Ito palette.}
\label{fig:recovery-s}
\end{figure*}

\emph{Exploratory}: Under rapid shocks (period~$=50$), V0+V1+V2+V3 entities
maintain higher mean alive counts (20.0) than V0-only (19.2), with a large
effect size (Cliff's~$\delta=0.69$, 95\%~CI $[0.60, 0.78]$).
Under slow shocks (period~$=200$), the effect persists ($\delta=0.66$,
$[0.56, 0.75]$), as shown in the centre panel of Figure~\ref{fig:recovery-s}.
The recovery time analysis shows that most conditions recover within a single
sampling interval (50 steps), suggesting that at the current temporal
resolution, recovery time differences are not reliably measurable.
This comparison is exploratory: shock resilience was not included in the
H1--H8 pre-registration.
The shock experiment uses only sparse harshness; extending to other harshness
levels would test whether the V3 buffering advantage is consistent across
resource regimes and is left as future work.

% ============================================================================
\section*{S4. Effect Size Robustness Check}
\label{sec:s4}

\emph{Exploratory}: The alive-count metric used in H1--H8 exhibits
ceiling/floor effects in some comparisons ($\delta=1.00$ with $[1.00,1.00]$~CI),
raising the concern that the metric is too coarse to distinguish
strong from overwhelming effects.
As a robustness check, we repeat the H1 and H2 comparisons using
\texttt{mean\_energy} at step~500---a continuous per-entity measure that
avoids the binary alive/dead dichotomy.
The H1 energy comparison at rich harshness yields $\delta=-0.61$
with CI $[-0.73, -0.47]$---\emph{reversed} relative to the alive-count
direction.
V1 entities have higher mean energy ($0.432 \pm 0.033$) than V0 entities
($0.402 \pm 0.017$), indicating that boundary protection benefits individual
energy budgets even though population size decreases.
The H2 energy comparison at rich yields $\delta=0.33$ with CI $[0.15, 0.49]$,
confirming that V3 benefits are detectable but not saturated on the continuous
metric.
In medium, sparse, and scarce, both comparisons remain at $\delta=1.00$.
These results demonstrate that the revised model produces a richer range of
effect sizes ($\delta$ from $-0.61$ to $1.00$) than the prior model's
universally saturated values.

% ============================================================================
\section*{S5. Pre-registered Sham Controls}
\label{sec:s5}

\emph{Exploratory report of pre-registered nulls}: the pre-registration states
that sham conditions should not differ from capability-absent counterparts.
Comparing V4 sham (\texttt{viroid\_v4\_sham}) against V3
(\texttt{viroid\_v0v1v2v3}) and V5 sham (\texttt{viroid\_v5\_sham}) against V4
(\texttt{viroid\_v0v1v2v3v4}) yields null effects at all harshness levels
($\delta=0.00$, $p=1.0$ for all eight comparisons).
These checks support the interpretation that the reported V4/V5 effects arise
from functional movement/lifecycle dynamics rather than merely running extra
computation.

% ============================================================================
\section*{S6. V4 Policy Evolution}
\label{sec:s6}

As an \textbf{exploratory} analysis (test seeds 100--199, no correction),
we extracted raw V4 policy weight vectors from alive V0+V1+V2+V3+V4 Viroid
entities at step~500 across $n\!=\!100$ test seeds.
The initial policy is $\mathbf{w}_0 = [0.5, 0.5, 0, 0, 0, 0, 0, 0]$
(equal gradient-following weights for $x$ and $y$; all others zero).
After 500 steps, mean weight deviations are small across all harshness levels
(maximum $|\Delta w| \approx 0.012$), with no weight exceeding $|\Delta w| = 0.02$.
The overall policy magnitude remains near the initial value ($0.71$--$0.73$
vs.\ initial $0.707$).
This constitutes a near-null result for policy evolution over 500 steps:
V4 chemotaxis behaviour is present and contributes to II$_B$, but the
weight vector has not diverged substantially from its initialisation.
This is consistent with the short simulation duration---500 steps may be
insufficient for evolutionary pressure to reshape the initial weight
distribution---and motivates longer runs in future work.

% ============================================================================
\section*{S7. Inter-Entity Competition}
\label{sec:s7}

As an \textbf{exploratory} analysis (seeds~0--99, no correction), we placed
Viroid~(V0 only) and Plasmid~(V0+V1+V2+V3) simultaneously in the same world,
using identical base parameters for both (an intentionally asymmetric capability
test: same physiology, different capability levels).
An ablation control placed Viroid~(V0) against Plasmid~(V0+V1+V2, no V3),
enabling attribution of any competitive advantage specifically to V3 metabolism.

In the primary condition (V3 present), Plasmid consistently dominates across
all harshness levels: frequency ratio 0.85~(rich), 0.74~(medium), 0.67~(sparse),
0.67~(scarce), with $\delta\!=\!1.0$ (primary vs.\ ablation, raw uncorrected
$p < 10^{-33}$).
In the ablation condition (V3 absent), Plasmid shows \emph{no competitive
advantage}: frequency ratio $\approx 0.43$--$0.50$, and Viroid slightly
dominates in resource-rich environments where V1+V2 overhead exceeds its benefit.
This directly attributes the competitive advantage to V3 metabolism specifically,
not to ``more capabilities'' in general.
The result confirms that V3 internal metabolism is not only the largest single
survival jump in the isolated ladder (H2) but also the decisive capability in
direct inter-entity competition.

% ============================================================================
\section*{S8. V1 Protection Regime Sweep}
\label{sec:s8}

Under the default damage parameters (composite hazard
$0.05 \times 0.05 = 0.0025$), V1 does not produce a net positive effect.
A supplementary V1 protection-regime sweep across 84 grid points
(7~damage probabilities $\times$ 4~damage amounts $\times$ 3~harshness levels,
seeds~0--39; data available upon acceptance)
reveals that \textbf{69 of 84 conditions show a net V1 survival
benefit} when environmental damage is present (measured by alive-count AUC).
The protection regime emerges above a composite hazard threshold of
$\sim$0.005 ($p_\text{damage} \times \text{damage\_amount}$): below this,
V1 cost exceeds its benefit; above, boundary maintenance becomes
protective.
The effect is strongest in medium, sparse, and scarce environments, where
V1 AUC advantages of 1{,}000--2{,}600 steps $\times$ entities are typical.
Additionally, V1 entities that survive show \emph{higher mean energy}
than V0-only: the boundary reduces energy leakage per step, meaning
each surviving entity is individually healthier.
This ``fewer but healthier'' pattern mirrors a $K$-strategy shift in
ecology---V1 reduces population quantity while increasing individual quality.

% ============================================================================

\footnotesize
\bibliographystyle{apalike}
\bibliography{references}

\end{document}
